\section{Zusammenfassung und Ausblick}

\subsection{Zusammenfassung und Beantwortung der Forschungsfrage}

Diese Arbeit untersuchte, wie zentrale Prinzipien von Attribute Stream-Based Access Control für Dateiöffnungen prototypisch mittels eBPF-LSM unter Linux umgesetzt werden können und welche technischen Einschränkungen sich insbesondere bei der nachträglichen Kontrolle bestehender Dateizugriffe ergeben. Hierzu wurde ein Prototyp entwickelt, der textuelle Policies und dynamische Attribute im Userspace verarbeitet, in eine feste Kernelrepräsentation übersetzt und über gepinnte eBPF-Maps bereitstellt. Neue Dateiöffnungen werden unmittelbar im LSM-Hook \texttt{file\_open} bewertet. Bereits geöffnete File Descriptors werden nach einer erfolgreichen Policy- oder Attributaktivierung sowie beim Erreichen einer entscheidungsrelevanten Zeitgrenze im Userspace erneut untersucht. Ergibt die Nachbewertung eine Verletzung, versucht der Userspace-PEP, den betroffenen Deskriptor im Zielprozess zu schließen.

Die Forschungsfrage lässt sich damit wie folgt beantworten: Eine prototypische Übertragung ausgewählter ASBAC-Prinzipien ist durch eine hybride Architektur möglich. eBPF-LSM bildet einen kernelgestützten Durchsetzungspfad für neue Dateiöffnungen, während Userspace-Komponenten die komplexere Policyverarbeitung, die Aktualisierung dynamischer Attribute und die Nachkontrolle bestehender Deskriptoren übernehmen. Feste Datentypen und begrenzte Kontrollflüsse machen die Entscheidungslogik für eBPF ausführbar. Zwei Map-Bänke und getrennte Generationszähler sorgen dafür, dass neue Policy- und Attributstände erst nach ihrer vollständigen Vorbereitung aktiviert werden. Die Aufteilung über Tail Calls begrenzt die Komplexität der einzelnen eBPF-Programme. Kernelspace-PEP und Userspace-PEP verwenden dieselbe Policysemantik, dieselben Objektidentitäten und die Real UID als definierte Subjektidentität.

Die Evaluation bestätigt die technische Machbarkeit für den untersuchten Funktionsumfang und das dokumentierte Zielsystem. Alle 59 Unit- und Komponententests sowie alle 17 privilegierten End-to-End-Szenarien bestanden. Dabei wurden unter anderem statische und dynamische Deny-Entscheidungen, die Combining-Regel \emph{deny-overrides}, konsistente Generationenwechsel, die kontrollierte Ablehnung ungültiger Zustände und der selektive Entzug bestehender File Descriptors beobachtet. Der Kernel-Verifier akzeptierte das eBPF-Programm, und der LSM-Hook wurde erfolgreich angehängt. Die UID-Semantik stimmte auch bei voneinander abweichenden Real- und Effective-UIDs mit der festgelegten Real-UID-Semantik überein.

Die Laufzeitmessungen zeigen zugleich, dass die zusätzliche Kontrolle nicht kostenfrei ist. Der Median einer gemessenen Dateiöffnung stieg in der verwendeten Testumgebung von 3,845~$\mu$s ohne Runtime auf 4,389~$\mu$s bei aktiver Runtime ohne Policy und auf 4,412~$\mu$s bei einer passenden Permit-Policy. Dies entspricht in diesem Versuchsaufbau einem medianen Mehraufwand von 14,1 beziehungsweise 14,7~Prozent. Policy- und Attributänderungen beeinflussten neue Entscheidungen im Median nach rund 102 beziehungsweise 103~ms. Vom Austausch einer Policydatei bis zum beobachteten \texttt{EBADF} eines zuvor geöffneten Deskriptors vergingen im Median 58,91~ms. Diese Werte charakterisieren die dokumentierte virtuelle Testumgebung; sie stellen weder eine Echtzeitgarantie noch einen isolierten Mikrobenchmark ausschließlich des LSM-Hooks dar.

Insgesamt erreicht der Prototyp damit das Ziel eines Machbarkeitsnachweises. Er zeigt, wie dynamische Entscheidungsgrundlagen in einen eBPF-LSM-Pfad einbezogen und bestehende Deskriptoren nach relevanten Zustandsänderungen erneut bewertet werden können. Er implementiert jedoch weder die vollständige Publish-Subscribe-Architektur von SAPL noch einen allgemeinen eigenständigen ASBAC-PDP. Ein geöffneter File Descriptor wird nicht als Subscription verwaltet, sondern bei einem Auslöser erneut aus dem zu diesem Zeitpunkt sichtbaren Systemzustand rekonstruiert. Die Bezeichnungen Kernelspace-PEP und Userspace-PEP beschreiben die jeweilige Durchsetzungsfunktion; beide Komponenten enthalten zugleich die für ihren Pfad erforderlichen Funktionen zur Policyentscheidung.

\subsection{Grenzen der Ergebnisse}

Die stärkste Durchsetzung bietet der Kernelspace-PEP für den eng definierten Zeitpunkt einer neuen Dateiöffnung. Die nachträgliche Behandlung eines bestehenden Zugriffs ist dagegen eine ereignisbasierte Best-Effort-Maßnahme. Das Durchlaufen von \path{/proc}, die erneute Policybewertung und der mittels \texttt{ptrace} ausgeführte \texttt{close}-Systemaufruf bilden keine atomare Operation. Zwischen Identifikation und Entzug kann der Zielprozess eine Deskriptornummer schließen oder neu belegen. Der nebenläufige Stresstest beobachtete zwar keinen Fehlentzug, kann Race-Freiheit aber nicht beweisen. Ein nicht erreichbarer Zielprozess kann nicht zum Schließen seines Deskriptors gezwungen werden.

Das Schließen eines File Descriptors beendet außerdem nicht jede Nutzung, die aus einer früheren Dateiöffnung hervorgegangen ist. Die Evaluation zeigte, dass duplizierte und vererbte Deskriptoren in den untersuchten Prozessen gefunden und geschlossen wurden, eine bereits bestehende \texttt{mmap}-Abbildung jedoch weiterhin lesbar blieb. Bereits abgeschlossene Lesevorgänge können ebenso wenig rückgängig gemacht werden. Der Prototyp stellt deshalb keine vollständige fortdauernde Vermittlung sämtlicher Dateinutzungen und keinen vollständigen Reference Monitor dar.

Weitere Grenzen folgen aus dem festgelegten Implementierungsumfang. Kontrolliert wird ausschließlich der LSM-Hook \texttt{file\_open}; andere Dateioperationen und Ressourcenarten werden nicht erfasst. Der Entzugsmechanismus ist an \texttt{x86\_64}-Linux und ausreichende \texttt{ptrace}-Berechtigungen gebunden. Policies identifizieren Ressourcen über Device und Inode, sodass der Austausch einer Datei durch ein neues Kernelobjekt eine erneute Policyverarbeitung erfordert. Die Real UID bildet effektive und Filesystem-UIDs sowie weitergehende Namespace-Semantiken nicht ab. Feste Map-Größen begrenzen pro Bank die Zahl statischer Policies und Streampolicies auf jeweils 16, die Zahl dynamischer Bedingungen pro Streampolicy auf vier und den gesamten Attributspeicher auf 1024 Einträge.

Auch die Evaluation besitzt begrenzte externe Validität. Die Versuche wurden auf einer einzelnen virtuellen Maschine, einem bestimmten Kernel und synthetischen Arbeitslasten ausgeführt. Die Latenzmessungen für Zustandsänderungen und Entzug umfassen jeweils nur zehn Wiederholungen; ihre p95-Werte sind daher keine belastbaren Schätzungen seltener Ausreißer. Ein Vergleich mit etablierten Linux-MAC-Systemen oder einer vollständigen SAPL-Laufzeit wurde nicht durchgeführt. Darüber hinaus war die definierte Gesamtprüfkette wegen eines verbleibenden Clippy-Strukturhinweises noch nicht vollständig erfolgreich, obwohl die Tests, die Formatprüfung und der Release-Build bestanden. Die Resultate belegen somit das Verhalten in den beschriebenen Szenarien, aber keine allgemeine Sicherheit, Skalierbarkeit oder Produktionstauglichkeit.

\subsection{Ausblick}

Ein vorrangiger Ansatzpunkt für weitere Arbeiten ist eine robustere Semantik für fortdauernde Zugriffe. Statt Nutzungssituationen nach einem Auslöser über \path{/proc} zu rekonstruieren, könnte eine Architektur Zugriffe und Entscheidungsänderungen ausdrücklich registrieren und als korrelierte Decision Streams verwalten. Damit ließe sich der Ansatz stärker an die Publish-Subscribe-Eigenschaften von ASBAC annähern. Zu untersuchen wäre, welche Ereignisse und Identifikatoren benötigt werden, um einen Zugriff über seinen Lebenszyklus stabil zu verfolgen, ohne einen File Descriptor selbst mit einer Subscription gleichzusetzen.

Für eine stärkere Durchsetzung müssten zusätzliche Zugriffspfade berücksichtigt werden. Weitere geeignete LSM-Hooks könnten Dateioperationen nach dem Öffnen vermitteln; ergänzende Mechanismen wären für Speicherabbildungen und bereits laufende E/A erforderlich. Besonders relevant ist eine Lösung, die das Zeitfenster zwischen der Identifikation eines Deskriptors und dessen Entzug verkleinert oder beseitigt. Dabei sind die Auswirkungen auf Nebenläufigkeit, Prozessstabilität, Berechtigungen und Portabilität systematisch zu untersuchen. Ein kooperativer Userspace-PEP im geschützten Prozess oder eine geeignete kernelgestützte Revocation-Schnittstelle könnten Alternativen zum derzeitigen \texttt{ptrace}-Eingriff bilden.

Die Policy- und Identitätsmodelle bieten ebenfalls Erweiterungspotenzial. Zusätzliche LSM-Hooks und Ressourcenarten erfordern eigene Anfrage- und Map-Strukturen. Effektive UID, Filesystem-UID, Capabilities, Cgroups sowie User- und Mount-Namespaces könnten als ausdrücklich typisierte Subjekt- oder Kontextattribute modelliert werden. Für größere Policybestände wären indexierte Suchstrukturen oder eine vorgelagerte Auswahl anwendbarer Policies zu untersuchen, da die derzeitigen festen linearen Scans bewusst auf kleine Mengen begrenzt sind.

Schließlich sollte eine Weiterentwicklung die Vertrauensgrenzen des Gesamtsystems härten. Dazu gehören minimale Capabilities für die Userspace-Komponenten, restriktive Rechte an Policy- und Attributquellen sowie an den gepinnten Maps, eine nachvollziehbare Herkunft und gegebenenfalls Signatur von Eingabedaten und ein individuelles Auditprotokoll für Entscheidungen und Entzugsversuche. Ergänzende Evaluationen auf physischer Hardware, weiteren Kernelversionen und unter realistischeren Mehrbenutzer- und Containerlasten könnten die Übertragbarkeit der Ergebnisse untersuchen. Größere Stichproben, isoliertere Hook-Mikrobenchmarks und Vergleiche mit bestehenden MAC- und ASBAC-Systemen würden zudem eine belastbarere Bewertung von Laufzeitaufwand und praktischem Nutzen ermöglichen.

\newpage
